%----------------------------------------------------------------------------------------
% METADATA BLOCK (YAML)  — see inbox/README.md
%----------------------------------------------------------------------------------------
% uid: ""
% school: "lycee"
% level: "3eme"
% section: "Economie et Gestion"
% chapter: "Statistiques"
% title: "Pointures vendues"
% difficulty: "Facile"
% tags: ["statistiques", "mode", "moyenne", "mediane", "frequences"]
% provenance: "own"
% batch: "2026-11-19-3eco-statistiques-n5"
%----------------------------------------------------------------------------------------

\newcommand{\R}{\mathbb{R}}
\newcommand{\vect}[1]{\overrightarrow{#1}}

% ---STATEMENT---
Un magasin de chaussures a relevé les pointures des $40$ paires vendues un samedi :
\begin{center}
\begin{tabular}{|l|c|c|c|c|c|c|}\hline
Pointure $x_i$ & $38$ & $39$ & $40$ & $41$ & $42$ & $43$ \\\hline
Effectif $n_i$ & $4$ & $6$ & $12$ & $10$ & $4$ & $4$ \\\hline
\end{tabular}\end{center}
\begin{enumerate}
  \item Donner le mode et l'étendue de la série.
  \item Calculer les fréquences en pourcentage. Quel pourcentage des paires vendues a une pointure au moins égale à $41$ ?
  \item Calculer la pointure moyenne.
  \item Dresser le tableau des effectifs cumulés croissants, puis déterminer la médiane.
\end{enumerate}

% ---SOLUTION---
\begin{enumerate}[itemsep=14pt]
  \item
  \begin{prop}[Mode, étendue]
  Le \textbf{mode} est la valeur de plus grand effectif ; l'\textbf{étendue} vaut \ib{$x_{\max}-x_{\min}$}.
  \end{prop}
  Le plus grand effectif est $12$ : le mode est $40$. \qquad Étendue : $43-38=5$.

  \item
  \begin{prop}[Fréquence]
  \ib{$f_i=\dfrac{n_i}{N}$}, que l'on multiplie par $100$ pour l'avoir en pourcentage.
  \end{prop}
  Ici $N=40$, donc $f_i=\dfrac{n_i}{40}\times100=2{,}5\,n_i\,\%$ :
  \begin{center}
  \begin{tabular}{|l|c|c|c|c|c|c|}\hline
  $x_i$ & $38$ & $39$ & $40$ & $41$ & $42$ & $43$ \\\hline
  $f_i$ (\%) & $10$ & $15$ & $30$ & $25$ & $10$ & $10$ \\\hline
  \end{tabular}\end{center}
  Pointure au moins égale à $41$ : $25+10+10=45\,\%$.

  \item
  \begin{prop}[Moyenne]
  \ib{$\overline x=\dfrac{n_1x_1+n_2x_2+\dots+n_px_p}{N}$}.
  \end{prop}
  \[\overline x=\dfrac{4\times38+6\times39+12\times40+10\times41+4\times42+4\times43}{40}\]
  \[\overline x=\dfrac{152+234+480+410+168+172}{40}=\dfrac{1\,616}{40}=40{,}4\]

  \item Effectifs cumulés croissants :
  \begin{center}
  \begin{tabular}{|l|c|c|c|c|c|c|}\hline
  $x_i$ & $38$ & $39$ & $40$ & $41$ & $42$ & $43$ \\\hline
  ECC & $4$ & $10$ & $22$ & $32$ & $36$ & $40$ \\\hline
  \end{tabular}\end{center}
  \begin{prop}[Médiane, effectif pair]
  Si $N$ est pair, la médiane est la moyenne des valeurs de rangs \ib{$\dfrac N2$ et $\dfrac N2+1$}.
  \end{prop}
  $N=40$ : on cherche les valeurs de rangs $20$ et $21$. Les ECC passent de $10$ à $22$ en $40$ : ces deux valeurs valent $40$.
  \[\text{Me}=\dfrac{40+40}{2}=40\]
\end{enumerate}
